\ssr{ПРИЛОЖЕНИЕ А}
\begin{center}\textbf{Краткое руководство по эксплуатации}\end{center}

Проект целиком размещается в целевой ОС Linux. Необходимы компилятор, Make, Python~3, заголовки работающего ядра и утилиты управления модулями. Физическое USB-устройство должно быть доступно целевой системе. Команды выполняются из каталога проекта:
\begin{lstlisting}[language=bash,basicstyle=\small\ttfamily]
make
modinfo -F vermagic ./xppen_mouse.ko
uname -r
sudo make load
sudo evtest
sudo make unload
\end{lstlisting}

Версия выпуска в \texttt{vermagic} должна совпадать с \texttt{uname -r}. Переменная \texttt{KDIR} автоматически указывает на заголовки работающего ядра. После обновления ядра модуль пересобирается.

В \texttt{evtest} выбирается созданное устройство относительной мыши. Касание удерживает левую кнопку, основная боковая кнопка --- правую. Выход из утилиты выполняется сочетанием Ctrl+C. Для обычного управления указателем диагностическая утилита не требуется.

При переподключении выполните \texttt{sudo make bind}. Чувствительность изменяется после выгрузки модуля. Пример загрузки с новым делителем: \texttt{sudo make load DIVISOR=4}. Допустимы значения делителя от 1 до 1024.

Команда \texttt{sudo make unload} выгружает модуль и восстанавливает предыдущие привязки. При ошибке восстановления после устранения причины команда повторяется. Состояние сохраняется в \path{/run/xppen-mouse/bindings.json}. Диагностика доступна через \texttt{sudo dmesg}. Очистка результатов сборки выполняется командой \texttt{make clean}.
\clearpage

\ssr{ПРИЛОЖЕНИЕ Б}
\begin{center}\textbf{Исходные тексты в электронном виде}\end{center}

Полный исходный текст USB-драйвера с комментариями включён в электронное приложение \path{appendix/main.c}. Файл является копией \path{src/main.c} из каталога проекта. Электронное приложение передаётся совместно с настоящей запиской.

Дополнительные исходные материалы:
\begin{itemize}
\item \path{appendix/Makefile} --- правила сборки, привязки, загрузки, выгрузки и очистки;
\item \path{appendix/manage.py} --- управление USB-привязками с сохранением состояния и откатом;
\item \path{appendix/clean.py} --- удаление артефактов сборки.
\end{itemize}

Оригиналы схем находятся в каталоге \path{diagram}. Подготовлены файлы:
\begin{itemize}
\item \path{probe.graphml} и \path{complete.graphml} --- подключение и завершение запроса;
\item \path{decode.graphml} и \path{stop.graphml} --- преобразование и остановка;
\item \path{recovery.graphml} и \path{binding.graphml} --- восстановление и переключение привязок.
\end{itemize}

Модель IDEF0 находится в каталоге \path{idef0}. Данные журнала, использованные в исследовании, представлены файлами \path{data/button_events.csv} и \path{data/motion_excerpt.csv}. Рисунок состояний кнопок создаётся сценарием \path{data/plot_events.py}, модель IDEF0 --- сценарием \path{idef0/go.py}.

Документ собирается из корня проекта сценарием \path{report_builder.sh}. Промежуточные файлы удаляются сценарием \path{clean.sh}.
\clearpage
